\studynote{4}{Wed 12 Oct 2022 18:59}{Concepts of Continuity}
\section{Continuity and Hölder continuity test}
\subsection{Definitions and Examples}

The main reference in this note is \cite{durettProbabilityTheory}.

\begin{definition}[Lipschitz continuous]
	$f$ is said to be \textit{Lipschitz continuous} if there is a constant $C$ so that $|f(x)-f(y)| \leq C \rho(x, y)$.
\end{definition}

Lipschitz continuity is a stronger notion of continuity than classical continuity. Lipschitz continuity implies continuity.

Geometrically, Lipchitz condition puts a finite bound on the slope of any secant line one can get from the graph of the function.

\begin{definition}
	A real or complex valued function $f$ on $d$-dimensional Euclidean space satisfies a Hölder condition with exponent $\alpha$, or is $\alpha$-Hölder continuous, when there are nonnegative real constants $C, a>0$, such that
\[
|f(x)-f(y)| \leq C\|x-y\|^\alpha
\]
for all $x$ and $y$ in the domain of $f$.
\end{definition}

\begin{itemize}
	\item When $\alpha > 1$, an $\alpha$-Hölder continuous function is \textit{constant}.
	\item When $\alpha = 1$, an $\alpha$-Hölder continuous function is \textit{Lipchitz continuous}.
	\item When $\alpha >0$, an $\alpha$-Hölder continuous function is \textit{uniformly continuous}.
	\item Whenver $0 < \alpha \le  \alpha'$, $\alpha'$-Hölder continuity implies $\alpha$-Hölder continuity.
\end{itemize}

\begin{remark}
	(from wikipedia) We have the following chain of strict inclusions for functions over a closed and bounded non-trivial interval of the real line:

Continuously differentiable $\subset $ Lipschitz continuous $\subset \alpha$-Hölder continuous $\subset$ uniformly continuous $\subset$ continuous

where $\alpha \in (0,1]$.
\end{remark}

\begin{theorem}
8.1.5. Brownian paths are Hölder continuous for any exponent $\gamma<1 / 2$.
\end{theorem}

\begin{theorem}
	8.1.6. With probability one, Brownian paths are not Lipschitz continuous (and hence not differentiable) at any point.
\end{theorem}

\subsection{Test for hölder continuity exponent}

From the definition of hölder continuity: 
\[
	|f(x)-f(y)| \leq C\|x-y\|^\alpha
\] 
Our goal is to find smallest possible $\alpha$ such that the equation above holds true empirically.

We take logarithm on both sides (using the fact that $\log$ is monotone increasing):
\[
	\log |f(x) - f(y)| \le \log C + \alpha \log \|x-y\|
\] 
rearrange and assuming $\log\|x-y\|<0$ for small $\|x-y\|$:
\begin{equation}
	\alpha \le \frac{\log \left| f(x) - f(y) \right| - \log C}{\log \|x - y\|} \sim \frac{\log \left| f(x) - f(y) \right| }{\log \|x-y\|}
.\end{equation} 
The smallest such $\alpha$ is then
\begin{equation}
	\label{eqn:holder-fit-1}
	\alpha 
	= \operatorname{inf}_{x, y} \frac{\log \left| f(x) - f(y) \right| - \log C}{\log \|x - y\|}
.\end{equation}


The term involving $\log C$ should fade away when $\|x-y\|$ is sufficiently small. 

\subsubsection{Compute $\alpha$ by finding a limit}
Let $x_1, \ldots, x_i, \ldots, x_n$ be random points in the domain of $f$. Let $\varepsilon_1, \ldots, \varepsilon_j, \ldots, \varepsilon_m$ be small changes.

Denote 
		\[
			\hat{a} _i^{(j)} := \frac{\log \left| f(x_i) - f(x_i + \varepsilon_j) \right| - \log C}{\log \varepsilon_j}
		.\] 
	We know $\hat{a}_i^{(j)}$ is always an \textbf{upper estimate} of $\alpha$. Thus a reasonable guess of $\alpha$ would be:
	\[
		\hat{\alpha} = \operatorname{inf}_{i, j} \hat{a}_i^{(j)}
	.\] 

Since $\hat{a}_i^{(j)} \sim a_i^{(j)}$ as $j \to  \infty $, we can replace our estimate of $\alpha$ by
	\[
		\hat{\alpha}  = \lim_{j \to \infty} \operatorname{inf}_i \hat{a} _i^{(j)}
	.\] 

	The limit, however, converges very slowly. As $\varepsilon_j$ decreases exponentially, $\left| \log \varepsilon_j \right| $ only increases linearly, so the error terms decreases very slowly.

\subsubsection{Find $\alpha$ and $C$ by Curve Fitting}

	Rewrite equation \eqref{eqn:holder-fit-1} as
	\[
		\alpha 
		= \operatorname{inf}_{x, \varepsilon} \frac{\log \left| f(x) - f(x+\varepsilon) \right| - \log C}{\log \varepsilon}
		\quad\text{where }\varepsilon>0
	.\]
Hence
\[
		\alpha 
		= \operatorname{inf}_{\varepsilon}
		\left( \operatorname{inf}_x
		\frac{\log \left| f(x) - f(x+\varepsilon) \right|}{\log \varepsilon}
-\frac{\log C}{\log \varepsilon} \right) 
	.\]

We claim:
\begin{proposition}
\[
		\alpha 
		=\lim_{\delta \to  0}
		\operatorname{inf}_{\varepsilon<\delta}
		\left( \operatorname{inf}_x
		\frac{\log \left| f(x) - f(x+\varepsilon) \right|}{\log \varepsilon}
-\frac{\log C}{\log \varepsilon} \right) 
	.\]
\end{proposition}
	This is basically saying that all roughness must remain in arbitrarily small detail.
\begin{proof}
	Fix some $\varepsilon'$ such that 
	for some $\delta>0$ and all $\varepsilon < \varepsilon'$, \[
		 \operatorname{inf}_x
		\frac{\log \left| f(x) - f(x+\varepsilon) \right|}{\log \varepsilon}
-\frac{\log C}{\log \varepsilon} 
> \alpha + \delta
	.\] 
	We claim that $f$ is actually $(\alpha + \delta)$-Hölder continuous, which results in a contradiction. It suffices to show that there exists some constant $C$ such that for all $\varepsilon$ where $\varepsilon\ge \varepsilon'$, \[
		\left| f(x) - f(x+\varepsilon) \right| \le C \varepsilon^{\alpha+\delta}
	.\] 
But since $f$ is $\alpha$-Hölder continuous, there exists some $C'>0$ such that
\[
		\left| f(x) - f(x+\varepsilon) \right| \le C' \varepsilon^{\alpha}
.\] 
Take $C = \frac{C'}{(\varepsilon')^{-\delta}}$, and we have
\[
		\left| f(x) - f(x+\varepsilon) \right| \le C (\varepsilon')^{\delta} \varepsilon^{\alpha}
		\le C \varepsilon^{\delta} \varepsilon^{\alpha}
		= C' \varepsilon^{\alpha+\delta}
.\]
\end{proof}
This claim allows us to take smaller and smaller $\varepsilon$, and be convinced that the limit infimum will be our desired $\alpha$.

%The next claim gives a reasonable error bound for dropping the infimum over $\varepsilon$ when calculating $\alpha$.
\begin{proposition}
	We have the bound
	\[
		 \operatorname{inf}_x
		\frac{\log \left| f(x) - f(x+\varepsilon) \right|}{\log \varepsilon}
-\frac{\log C}{\log \varepsilon}  
\le 
		\operatorname{inf}_{\varepsilon}
		\left( \operatorname{inf}_x
		\frac{\log \left| f(x) - f(x+\varepsilon) \right|}{\log \varepsilon}
-\frac{\log C}{\log \varepsilon} \right) 
+ O(\varepsilon)
	\]
	under the following assumptions:

\begin{enumerate}
	\item $f$ has bounded total variation.
\end{enumerate}

\color{gray}Note: I don't know how to prove this theorem. I discussed this with Zijie, and have some sketchy idea on how to prove it for brownian motions. 
\end{proposition}

Using the approximation above, for some specific $\varepsilon$, we can approximate $\alpha$ by
\[
	\alpha \approx \operatorname{inf}_x
		\frac{\log \left| f(x) - f(x+\varepsilon) \right|}{\log \varepsilon}
-\frac{\log C}{\log \varepsilon} 
.\] 

We then use $\operatorname{inf} _i a_i^{(j)}$ as an upper estimation:

\[
\operatorname{inf} _i a_i^{(j)} \approx
	 \operatorname{inf}_x
		\frac{\log \left| f(x) - f(x+\varepsilon) \right|}{\log \varepsilon}
.\] 

We finally have a curve to fit:
\begin{equation}
	\label{eqn:holder-fit-curve}
	\alpha + \frac{\log C}{\log \varepsilon_j} = \operatorname{inf} _i a_i^{(j)} + O(\varepsilon) 
.\end{equation} 

{\color{gray} Note: We still need to bound the errors of the approximations above.}

Thus our testing strategy is:
\begin{enumerate}
	\item Pick some small $\varepsilon_1$ and randomly pick $x_1^{(1)}, \ldots, x_n^{(1)}$ in the domain of interested function.
	\item Calculate
		\[
			a_i^{(j)} := \frac{\log \left| f(x_i^{(j)}) - f(x_i^{(j)} + \varepsilon_j) \right| }{\log \varepsilon_j}
		.\] 

	\item Calculate $\operatorname{inf} _i a_{i}^{(j)}$
	\item Do the same for successively smaller $\varepsilon_2, \ldots, \varepsilon_m, \ldots$ that approach $0$.
\item Fit the curve \eqref{eqn:holder-fit-curve}. Then read the $\alpha$ and $C$.
\end{enumerate}


